\documentclass[11pt,a4paper]{article}
\usepackage[T1]{fontenc}
\usepackage[utf8]{inputenc}
\usepackage[french]{babel}
\usepackage{lmodern}
\usepackage{amsmath,amssymb}
\usepackage[margin=2.35cm]{geometry}
\usepackage{booktabs,longtable,tabularx}
\usepackage{enumitem}
\usepackage{graphicx}
\usepackage{xcolor}
\usepackage{microtype}
\usepackage{hyperref}

\hypersetup{colorlinks=true,linkcolor=blue!45!black,urlcolor=blue!45!black}
\setlist{nosep}
\newcommand{\code}[1]{\texttt{#1}}
\newcommand{\Kaut}{K^{*}_{\mathrm{aut}}}
\newcommand{\fait}{\textbf{[Fait observé]}}
\newcommand{\inference}{\textbf{[Inférence]}}
\newcommand{\hyp}{\textbf{[Hypothèse]}}
\newcommand{\incertitude}{\textbf{[Incertitude]}}

\title{M4.2B --- rapport intermédiaire\\
\large Lecture de la campagne d'exploration (87 runs) pendant que la
confirmation tourne}
\author{Programme M4.2B --- dossier \code{m4\_2b\_credit\_soc}}
\date{5 août 2026}

\begin{document}
\maketitle

\begin{abstract}
Ce document lit exclusivement les données de la campagne
d'\textbf{exploration} (87 runs). Une campagne de
\textbf{confirmation}, plus petite mais statistiquement plus exigeante,
est en cours d'exécution au moment de la rédaction et n'est pas encore
intégrée ici. \textbf{Rien dans ce document n'a été validé par cette
confirmation} --- c'est une lecture provisoire, à seule fin de savoir
où en est la science en attendant. Le document est autonome : les
notions techniques utilisées (K0, $\gamma$, $\eta$, $\delta$, $\sigma$,
exposants de queue, critères de confirmation, etc.) sont définies dans
les sections~\ref{sec:reperes} et~\ref{sec:criteres} avant d'être
utilisées.
\end{abstract}

\tableofcontents

\section{Contexte : le programme M4.2B en deux phrases}

M4.2B est une simulation multi-agents d'une économie de crédit : des
entités abstraites empruntent, prêtent, se versent des intérêts, et
peuvent faire faillite --- une faillite pouvant en déclencher d'autres
en cascade (« avalanche »). Le programme poursuit deux objectifs
\textbf{simultanés} :

\begin{itemize}
\item \textbf{Objectif A} --- Le revenu d'intérêt réellement perçu par
  chaque entité à un instant donné a-t-il, dans une région significative
  de l'espace des paramètres, une distribution à queue épaisse de type
  \textbf{Pareto} (quelques entités touchant un revenu disproportionné
  selon une loi de puissance), et cette queue peut-elle être pilotée de
  façon reproductible en agissant sur la structure du marché du
  crédit --- en particulier sur l'intensité des appariements
  emprunteur/prêteur (paramètre $\eta$, voir \S\ref{sec:reperes}) ?
\item \textbf{Objectif B} --- Quel que soit le réglage trouvé pour
  l'objectif A, il ne doit pas détruire une propriété déjà établie dans
  les versions antérieures du modèle : la taille des cascades de
  faillites suit une loi de puissance, signature d'une dynamique proche
  d'un point critique. Un réglage qui produit une belle queue Pareto des
  revenus mais éteint cette criticalité n'est pas un succès --- c'est un
  compromis à documenter comme tel.
\end{itemize}

\section{Méthode : exploration puis confirmation}

Une \textbf{campagne d'exploration} balaie 29 combinaisons de
paramètres (« cellules »), chacune répétée sur 3 graines aléatoires (0,
1, 2) pour vérifier une cohérence minimale, sur 3000 pas de temps
simulés (une cellule est poussée à 10\,000 pas pour vérifier la
stabilité à long terme) --- 87 runs au total, tous terminés avec
succès. Son rôle est de repérer où, dans l'espace des paramètres, il se
passe quelque chose d'intéressant, avant d'investir plus de calcul sur
un petit nombre de cellules.

Avant de lancer cette exploration, des \textbf{critères numériques
précis} ont été fixés à l'avance (avant de voir le moindre résultat)
pour décider si un effet observé peut être qualifié de « robuste » ---
un garde-fou classique contre le biais qui consiste à ajuster une
conclusion après coup pour qu'elle colle aux données. Une
\textbf{campagne de confirmation}, en cours au moment de la rédaction,
applique ces critères sur un petit nombre de cellules choisies dans
l'exploration, avec 5 graines \textbf{nouvelles} (10 à 14, jamais vues
pendant l'exploration) --- précisément pour empêcher de choisir a
posteriori les graines qui arrangent le résultat.

\section{Repères : ce que mesurent les grandeurs citées plus loin}
\label{sec:reperes}

\begin{itemize}
\item \textbf{$K_0$} : le capital de départ d'une entité nouvellement créée.
\item \textbf{$\Kaut$} : le capital vers lequel une entité isolée (sans
  jamais emprunter ni prêter) convergerait naturellement à long terme,
  calculé à partir des équations de production/dépréciation du modèle.
  C'est une échelle de référence : ce qui compte pour la dynamique
  n'est pas la valeur brute de $K_0$, mais son rapport à $\Kaut$
  (est-on loin ou proche de l'échelle naturelle de l'entité ?).
\item \textbf{$\gamma$} : la concavité de la fonction de production ---
  à quelle vitesse les rendements du capital diminuent. $\gamma$ petit
  = rendements fortement décroissants ; $\gamma$ proche de 1 =
  rendements presque linéaires.
\item \textbf{$\eta$} : le nombre de tentatives d'appariement
  emprunteur/prêteur à chaque pas de temps --- l'intensité d'activité
  du marché du crédit. Deux familles sont testées : une famille
  \textbf{linéaire} (le nombre de tentatives est proportionnel à la
  population active, avec un facteur $\rho$ ; $\rho=1$ est la référence
  historique) et une famille \textbf{non linéaire} (exposant $\beta$
  autour d'une population de référence).
\item \textbf{$\delta$, $\sigma$} : $\delta$ est le taux de
  dépréciation du capital à chaque pas ; $\sigma$ est l'amplitude des
  chocs aléatoires individuels sur la production. Ensemble, ils fixent
  la « vitesse » et le « bruit » de l'économie.
\item \textbf{Mécanisme \emph{deg\_out}/\emph{rq}} (établi lors d'une
  phase antérieure) : la dispersion du revenu d'intérêt entre entités
  se décompose en une part due au \textbf{nombre de contrats de prêt
  actifs accumulés} (\emph{deg\_out}, un effet de quantité) et une part
  due au \textbf{taux $\times$ montant remboursé par contrat}
  (\emph{rq}, un effet de prix). En baseline, environ 77\,\% de la
  dispersion vient du nombre de contrats, 23\,\% du prix par contrat.
\item \textbf{$\hat\alpha$ et stabilité de seuil} : $\hat\alpha$ est
  l'exposant ajusté d'une loi de puissance sur la queue de la
  distribution des revenus. Un ajustement de queue nécessite de choisir
  un seuil au-delà duquel les données sont considérées comme « la
  queue » ; la \textbf{stabilité de seuil} consiste à refaire
  l'ajustement avec un seuil deux fois plus bas et à vérifier que
  $\hat\alpha$ ne bouge pas de plus de 20\,\% --- sinon, la « queue
  Pareto » n'est qu'un artefact du choix de seuil, pas un vrai trait des
  données.
\item \textbf{$n_{\mathrm{tail}}$} : le nombre de points de données
  effectivement utilisés pour l'ajustement de la queue --- trop peu de
  points rend l'exposant peu fiable.
\item \textbf{Test de Vuong} : un test statistique qui compare deux
  modèles de distribution concurrents (ici, loi de puissance contre
  exponentielle, ou loi de puissance contre lognormale) et indique
  lequel est mieux supporté par les données.
\item \textbf{$\hat\tau$ et rapport de branchement} : $\hat\tau$ est
  l'exposant de loi de puissance ajusté sur la taille des avalanches de
  faillites. Le rapport de branchement est le nombre moyen de faillites
  supplémentaires déclenchées par une faillite donnée --- plus il est
  proche de 1, plus la dynamique de cascade est proche d'un régime
  critique auto-entretenu.
\item \textbf{Plancher de renouvellement} : on repère les entités dans
  le décile le plus riche (top 10\,\% en capital ou en revenu) au début
  d'une fenêtre d'observation, puis on mesure quelle fraction de ce
  groupe initial est \emph{encore} dans le top 10\,\% en fin de
  fenêtre. Un plancher proche de 0 signale un renouvellement total de
  l'élite (grande mobilité) ; proche de 1, une élite figée.
\end{itemize}

\section{Les critères de confirmation, fixés avant tout résultat}
\label{sec:criteres}

Une queue Pareto des revenus d'intérêt sera qualifiée de robuste pour
une cellule si, sur au moins 3 des 5 graines de confirmation :

\begin{enumerate}
\item \textbf{Stabilité de seuil} (\S\ref{sec:reperes}) : écart de
  moins de 20\,\% relatif entre $\hat\alpha$ au seuil optimal et
  $\hat\alpha$ à seuil moitié.
\item \textbf{Au moins 100 points} dans la queue en moyenne par
  instantané.
\item \textbf{Signe et ampleur de l'effet cohérents} entre les graines,
  pas seulement en moyenne.
\item \textbf{Le test de Vuong favorise la loi de puissance} dans la
  majorité des instantanés (contre l'exponentielle et contre la
  lognormale).
\item \textbf{L'effet ne se réduit pas} à un simple changement
  d'inégalité globale (indice de Gini) ou au seul mécanisme de quantité
  de contrats (\emph{deg\_out}) --- il doit s'agir d'un vrai changement
  de forme de la queue.
\end{enumerate}

Un exposant d'avalanche sera qualifié d'indépendant de la taille du
système si : l'ajustement détecte une vraie coupure finie (pas de
saturation contre une borne numérique de l'algorithme) sur toutes les
tailles testées ; $\hat\tau$ varie de moins de 15\,\% relatif entre les
tailles de population les plus petites et les plus grandes ; le rapport
de branchement reste stable à $\pm 0{,}05$ près sur la même plage.

\textbf{Aucun chiffre de ce document n'a encore été confronté à ces
critères} --- c'est l'objet de la campagne de confirmation en cours.

\section{Le résultat le plus robuste de l'exploration : l'objectif A
échoue presque partout}

La fraction d'instantanés satisfaisant le critère de stabilité de seuil
(critère 1, \S\ref{sec:criteres}) est \textbf{quasi nulle sur 28 des 29
cellules} du grid --- autrement dit, la queue Pareto n'est stable au
changement de seuil dans (quasi) aucun instantané, quelle que soit la
cellule testée.

\fait{} C'est le résultat le plus solidement établi de toute la
campagne --- une régularité sur 87 runs couvrant $K_0$ (facteur 2000),
$\gamma$ (brut et compensé, voir \S\ref{sec:gamma}), $\eta$ linéaire
(facteur 32), $\eta$ non linéaire, $\delta$ et $\sigma$ conjoints, et un
réglage institutionnel de comparaison.

Seule exception : la cellule à $\delta=\sigma=0{,}10$ atteint
9{,}2\,\% --- toujours loin d'une majorité d'instantanés stables, mais
un ordre de grandeur au-dessus de tout le reste du grid (détail au
\S\ref{sec:deltasigma}).

\section{$K_0$ : le mécanisme \emph{deg\_out} se confirme, mais pas de
façon monotone simple}

\begin{table}[h]
\centering
\begin{tabular}{@{}rrrrrr@{}}
\toprule
$K_0$ & $\hat\alpha$ & $n_{\mathrm{tail}}$/inst. & branchement & pop.\ finale & plancher renouv. \\
\midrule
1    & 3{,}59 & 105  & 0{,}588 & 387  & \textbf{0{,}000} \\
5    & 4{,}08 & 119  & 0{,}709 & 620  & 0{,}022 \\
25 (baseline) & 3{,}78 & 252 & 0{,}786 & 1147 & 0{,}088 \\
100  & 3{,}51 & 500  & 0{,}795 & 1978 & 0{,}201 \\
500  & 3{,}10 & 1330 & 0{,}746 & 3852 & 0{,}509 \\
2000 & \textbf{2{,}79} & 3314 & 0{,}655 & 7247 & \textbf{0{,}796} \\
\bottomrule
\end{tabular}
\caption{Branche $K_0$, moyenne sur 3 graines d'exploration.}
\end{table}

$\hat\alpha$ décroît globalement avec $K_0$ (queue plus lourde à $K_0$
élevé, cohérent avec le mécanisme \emph{deg\_out} : plus de population,
plus de contrats accumulés possibles), mais pas monotoniquement (bosse
à $K_0=5$). Le rapport de branchement est \textbf{en cloche} (pic à
$K_0=100$), pas monotone non plus --- une nuance par rapport à l'idée
d'un effet $K_0$ simple sur la criticalité des avalanches.

Le changement le plus spectaculaire est ailleurs : le \textbf{plancher
de renouvellement} passe de $0{,}000$ ($K_0=1$ : l'élite du décile
supérieur se renouvelle entièrement) à $0{,}796$ ($K_0=2000$ : l'élite
est quasi figée). \inference{} $K_0$ contrôle autant la mobilité
sociale de long terme que la forme de la queue instantanée des revenus
--- deux effets à ne pas confondre.

$K_0=1$ et $K_0=2000$ sont les deux extrêmes retenus pour la
confirmation.

\section{$\gamma$ : l'effet réel n'apparaît qu'une fois l'échelle
contrôlée}
\label{sec:gamma}

\begin{table}[h]
\centering
\begin{tabular}{@{}rrr@{}}
\toprule
$\gamma$ & $\hat\alpha$ ($K_0=25$ fixe) & $\hat\alpha$ ($K_0$ compensé pour tenir $K_0/\Kaut$ constant) \\
\midrule
1/3          & 3{,}80 & \textbf{4{,}86} \\
0{,}4        & 3{,}81 & \textbf{4{,}34} \\
0{,}5 (baseline) & 3{,}78 & 3{,}78 \\
0{,}6        & 3{,}74 & \textbf{3{,}48} \\
2/3          & 3{,}35 & \textbf{3{,}31} \\
\bottomrule
\end{tabular}
\caption{Branche $\gamma$ : $K_0$ fixe vs.\ $K_0$ compensé.}
\end{table}

À $K_0$ fixe, la branche est quasi plate : $K_0=25$ représente une
fraction très différente de $\Kaut(\gamma)$ selon $\gamma$ (985 à
$\gamma=1/3$ contre 970\,299 à $\gamma=2/3$), ce qui confond l'effet
d'échelle (distance à $\Kaut$) avec l'effet de courbure proprement dit.
Une fois $K_0$ ajusté pour tenir ce rapport constant, la décroissance
de $\hat\alpha$ avec $\gamma$ devient nette et quasi monotone
($4{,}86 \to 3{,}31$).

\inference{} (exploration seule) plus la production est concave
($\gamma$ petit), plus la queue des revenus est légère ; plus elle est
proche de la linéarité ($\gamma$ grand), plus elle s'alourdit. C'est le
signal le plus net et le plus propre de toute la campagne pour un
paramètre autre que $\eta$ --- mais la branche « $K_0$ compensé »
n'était pas dans la sélection initiale de confirmation (voir
\S\ref{sec:limite}), ajoutée depuis en complément.

\section{$\eta$ : le levier le plus systématique du grid --- linéaire
fort, non linéaire nul}

Famille linéaire ($\eta$ proportionnel à la population, facteur
$\rho$) : effet le plus net sur tous les diagnostics mesurés, dans le
sens attendu.

\begin{table}[h]
\centering
\begin{tabular}{@{}rrrrr@{}}
\toprule
$\rho$ & $\hat\alpha$ & $\hat\tau$ & branchement & Vuong vs.\ exp.\ (frac.\ favorisant Pareto) \\
\midrule
0{,}125 & 4{,}20 & 1{,}84 & 0{,}480 & 0{,}10 \\
0{,}25  & 4{,}09 & 1{,}89 & 0{,}584 & 0{,}57 \\
1 (baseline) & 3{,}78 & 1{,}41 & 0{,}786 & 1{,}00 \\
2      & 3{,}93 & 1{,}17 & 0{,}854 & 1{,}00 \\
4      & 4{,}05 & 0{,}99 & 0{,}891 & 1{,}00 \\
\bottomrule
\end{tabular}
\caption{Branche $\rho$, moyenne sur 3 graines d'exploration.}
\end{table}

\fait{} À $\rho=0{,}125$ et $0{,}25$, le test de Vuong favorise
l'\textbf{exponentielle} sur la loi de puissance dans la majorité des
instantanés (seulement 10\,\%/57\,\% en faveur de Pareto) : un marché
de crédit clairsemé ne produit quasiment jamais une queue
Pareto-compatible, même ponctuellement. Les deux extrêmes
($\rho=0{,}125$ et $\rho=4$) sont dans les cellules de confirmation.

Famille non linéaire (exposant $\beta$ entre 0{,}5 et 1{,}5 autour
d'une population de référence) : \textbf{aucun effet mesurable} sur
$\hat\alpha$ (3{,}78--3{,}84), $\hat\tau$ (1{,}42--1{,}44), rapport de
branchement (0{,}785--0{,}789) ou population finale (1070--1170).
Résultat négatif propre, mais valable uniquement sur la plage testée :
rien n'indique ce qui se passerait en dehors de $\beta \in
[0{,}5\,;\,1{,}5]$.

\section{$\delta$ et $\sigma$ conjoints : le seul endroit où l'objectif
A progresse un peu --- au prix de l'objectif B}
\label{sec:deltasigma}

\begin{table}[h]
\centering
\begin{tabular}{@{}lrrrr@{}}
\toprule
$(\delta,\sigma)$ & $\hat\alpha$ & stabilité de seuil (frac.) & branchement & $\hat\tau$ \\
\midrule
$(0{,}01\,;0{,}01)$ baseline & 3{,}78 & 0{,}004 & 0{,}786 & 1{,}41 \\
$(0{,}02\,;0{,}02)$ & 3{,}55 & 0{,}011 & 0{,}709 & 1{,}52 \\
$(0{,}05\,;0{,}05)$ & 3{,}19 & 0{,}011 & 0{,}547 & 1{,}71 \\
$(0{,}10\,;0{,}10)$ & 3{,}04 & \textbf{0{,}092} & \textbf{0{,}385} & 1{,}82 \\
$(0{,}05\,;0{,}25)$ --- régime historique bruité & 4{,}40 & 0{,}000 & \textbf{0{,}357} & 1{,}97 \\
\bottomrule
\end{tabular}
\caption{Branche $\delta/\sigma$ joints, moyenne sur 3 graines
d'exploration.}
\end{table}

À mesure que $\delta$ et $\sigma$ augmentent ensemble, le rapport de
branchement \textbf{s'effondre} ($0{,}786 \to 0{,}385$, plus de moitié
perdu) et $\hat\tau$ s'alourdit (avalanches moins critiques). Dans le
même mouvement, la stabilité de seuil progresse légèrement --- c'est la
seule branche du grid où elle bouge de façon notable.

\hyp{} (exploration seule) compromis direct entre les deux objectifs
sur cette branche --- davantage de bruit individuel semble aider
marginalement à stabiliser la queue des revenus, mais au prix d'une
nette dégradation de la criticalité des avalanches. Non confirmé au
moment de la rédaction ; cette cellule a été ajoutée en complément à la
confirmation (voir \S\ref{sec:limite}).

\section{Une limite identifiée dans le choix des cellules de
confirmation}
\label{sec:limite}

Trois façons de choisir les cellules à confirmer étaient possibles au
sortir de l'exploration : la taille de l'effet sur le mécanisme de prix
(\emph{rq}, \S\ref{sec:reperes}), la stabilité de seuil elle-même, ou
un comportement d'avalanche particulièrement intéressant. La sélection
initiale n'a utilisé que le premier critère. Deux cellules ressortaient
nettement sur les deux autres et n'étaient pas dans les 6 retenues :

\begin{itemize}
\item la cellule $\delta=\sigma=0{,}10$ --- seule cellule à sortir du
  bruit sur la stabilité de seuil elle-même (\S\ref{sec:deltasigma}) ;
\item la branche $\gamma$ compensée --- effet le plus propre sur
  $\gamma$ une fois l'échelle contrôlée (\S\ref{sec:gamma}).
\end{itemize}

Décision prise : ne pas modifier les 6 cellules déjà lancées en
confirmation (les changer après coup reviendrait à sélectionner a
posteriori les résultats qui arrangent), mais lancer ces cellules
supplémentaires \textbf{en complément}, sur un lot de calcul séparé,
décidé avant tout résultat de confirmation --- à rapporter séparément
dans le rapport final pour ne pas mélanger les deux sélections.

\section{Ce qui reste strictement en attente de la confirmation}

\begin{itemize}
\item Le critère 3 (\S\ref{sec:criteres}) --- cohérence du signe et de
  l'ampleur sur au moins 3 graines \textbf{différentes de celles de
  l'exploration} --- ne peut par construction pas être évalué avec les
  seules données d'exploration : c'est l'objet même de la confirmation.
\item Le critère 2 (\S\ref{sec:criteres}, au moins 100 points de queue
  en moyenne) est globalement satisfait, mais tout juste sur $K_0=1$
  (105) et sur la cellule $\gamma=2/3$ (104), deux des cellules en
  confirmation --- à surveiller sur les nouvelles graines.
\item Le critère 5 (\S\ref{sec:criteres}, l'effet ne se réduit pas à
  l'inégalité globale ou au seul mécanisme de quantité de contrats) n'a
  pas encore été vérifié systématiquement sur les cellules retenues ; à
  faire une fois la confirmation terminée.
\end{itemize}

\section*{À ne pas citer comme conclusion}

Aucun chiffre de ce document n'a passé les cinq critères de robustesse
du \S\ref{sec:criteres}. Le statut correct de toute affirmation
ci-dessus est « observé sur les graines d'exploration, non confirmé ».

\end{document}
